% New checkpoint34 source. Actual full-GT data and original signed marking.
\statement{A.111}{Proved}{A coherent factor-frame obstruction whose signed image is the original CM character}
Retain the actual full-GT negative involution \(\alpha=(-1,\phi)\),
the canonical norm word \(u\), the even dihedral exponent \(k\),
and \(d=1+2k\in\widehat{\mathbf Z}^{\times}\) of A.103--A.108.
Let \(N_1=\overline{[\ker q,F]}\), \(E_1=F/N_1\), and
\(B_1=B/N_1\). All actions below are genuine continuous actions
on these actual profinite quotients. Put
\[
 e=d\bmod3\in\{1,-1\},\qquad
 \varepsilon=(1-e)/2,\qquad
 R_e=y^{\varepsilon}r,\qquad v=u^{-1}s=su.
\tag{A.111.1}
\]
Here the two values of \(e\) are ordinary integer representatives.
Thus \(R_1=r\) and \(R_{-1}=sr^{-1}s\).
The automorphism \(\Psi_1\) of A.107 has
\(\Psi_1(s)=v\) and, by A.108, admits a pure frame
\[
 \Psi_1(r)=hR_eh^{-1},\qquad h\in E_1.
\tag{A.111.2}
\]
Use the ORIGINAL unique full-group frame \(A_r\) of A.105, and
use the following full-group words, then their images in \(B_1\):
\[
 J_e^{\alpha}=
 \begin{cases}s\phi A_r,&\varepsilon=0,\\
 A_rs,&\varepsilon=1.
 \end{cases}
\tag{A.111.3}
\]
Define the closed profinite group and its involution
\[
 T_e=C_{E_1}(R_e),\qquad
 \kappa_e(t)=(J_e^{\alpha})^{-1}\alpha(t)J_e^{\alpha}.
\tag{A.111.4}
\]
Then
\[
 w_h=(J_e^{\alpha})^{-1}\alpha(h)^{-1}vh
       \in Z^1(C_2,T_e),\qquad
 \omega_{\alpha}=[w_h]\in H^1(C_2,T_e)
\tag{A.111.5}
\]
is independent of the choice of the pure quotient frame \(h\).
We use continuous nonabelian cohomology, with cocycles
\(w\kappa_e(w)=1\) and equivalence
\(w\sim a^{-1}w\kappa_e(a)\), \(a\in T_e\).
The class \(\omega_{\alpha}\) is the distinguished point if and
only if a frame in (A.111.2) can be chosen to satisfy
\[
                  \alpha(h)J_e^{\alpha}=vh.
\tag{A.111.6}
\]
Furthermore the actual signed reduction restricts to a continuous
\(C_2\)-equivariant homomorphism
\[
 \chi_e=q_3|_{T_e}:T_e\longrightarrow\{I,-I\},
\tag{A.111.7}
\]
where the target has trivial action. Its induced map on pointed
cohomology sets satisfies
\[
 \boxed{\quad
 \chi_{e,*}(\omega_{\alpha})
       =\epsilon_A eI=(-1)^{\Delta(\alpha)}I.
 \quad}
\tag{A.111.8}
\]
The sign \(\epsilon_A\) and the character \(\Delta\) are exactly
those of A.105 and original LM7\_PUSH~E.65. Thus a nonzero original
CM character obstructs (A.111.6). This theorem neither decides
that character nor asserts that its vanishing suffices to make
\(\omega_{\alpha}\) the distinguished point.

\begin{proof}
\emph{A.111.1. The two actual negative conjugators.}
The original extensions to \(B\) satisfy
\(\alpha(s)=s\phi\), \(\alpha(sr)=(sr)^{-1}\), and
\(\Sigma=\alpha c_B\) has
\(\Sigma(r)=A_rrA_r^{-1}=\phi^{-1}r\).
Consequently the first word in (A.111.3) satisfies
\(\alpha(r)=J_1^{\alpha}r^{-1}(J_1^{\alpha})^{-1}\),
as already recorded in A.105. Since \(\alpha(s)=s\phi\)
is an involution, \((s\phi)^2=1\). Direct multiplication gives
\[
 \alpha(sr^{-1}s)=A_rrA_r^{-1}
      =(A_rs)(sr^{-1}s)^{-1}(A_rs)^{-1}.
\]
Therefore, for both branches,
\[
 \alpha(R_e)=J_e^{\alpha}R_e^{-1}(J_e^{\alpha})^{-1},
 \qquad \alpha(J_e^{\alpha})J_e^{\alpha}=1.
\tag{A.111.9}
\]
For the second equality we first work in the full group \(B\).
Applying \(\alpha\) twice to the first equality shows that
\(\alpha(J_e^{\alpha})J_e^{\alpha}\) centralizes \(R_e\).
Its image in \(S_3\) is one. The finite-factor centralizer
theorem \cite[Theorem B']{A_HR1985}, applied to \(r\) or its
displayed conjugate, gives \(C_B(R_e)\cap F=1\).
This proves the equality in the full group, hence in \(B_1\).
It does not invoke uniqueness of quotient frames.

The first equality in (A.111.9) shows that \(\kappa_e\) preserves
\(T_e\); the second shows that \(\kappa_e^2=1\).
Continuity follows from continuity of the group operations and
of the descended actual automorphism \(\alpha\).

\emph{A.111.2. The cocycle and its exact frame ambiguity.}
The intertwining of A.107 is
\(\alpha\Psi_1=\Psi_1c_B\), with \(\alpha(v)=v\)
and \(v^2=1\). Evaluate it on \(r\), using
\(c_B(r)=sr^{-1}s\), to obtain
\[
 (\alpha(h)J_e^{\alpha})R_e^{-1}
       (\alpha(h)J_e^{\alpha})^{-1}
       =(vh)R_e^{-1}(vh)^{-1}.
\tag{A.111.10}
\]
Thus \(w_h\) centralizes \(R_e\). Both
\(\alpha(h)J_e^{\alpha}\) and \(vh\) have the same
transposition as image in \(S_3\), so \(w_h\in E_1\).
This proves \(w_h\in T_e\). A literal calculation, in the
stated order, gives
\[
 \kappa_e(w_h)=h^{-1}v\alpha(h)J_e^{\alpha},
 \qquad \kappa_e(w_h)w_h=1.
\tag{A.111.11}
\]
The reversed norm \(w_h\kappa_e(w_h)=1\) follows too.

Every other pure frame in (A.111.2) is uniquely \(ht\) for
some \(t\in T_e\), because equality of the two conjugates
is exactly \(h^{-1}h'\in C_{E_1}(R_e)\). Substitution gives
\[
              w_{ht}=\kappa_e(t)^{-1}w_ht.
\tag{A.111.12}
\]
This is the stated cohomology equivalence with
\(a=\kappa_e(t)\), and every such equivalence occurs as
\(t\) varies. The class is therefore independent of \(h\).
It is the distinguished point precisely when some \(w_{ht}=1\),
which is (A.111.6). If a topology is desired, the cocycles form a
closed subspace of the compact group \(T_e\); their quotient by
this continuous compact-group action is compact Hausdorff.

\emph{A.111.3. The actual comparison to the signed scalar group.}
Since \(N_1\subseteq F'\), we have
\((E_1)_{\rm ab}=F_{\rm ab}=\widehat{\mathbf Z}^2\).
Conjugation by either \(R_e\) has on this abelianization the
same order-three matrix as \(r\), namely
\[
                  M=\begin{pmatrix}0&-1\\1&-1\end{pmatrix}.
\]
The determinant of \(M-I\) is three. Its adjugate and
injectivity of multiplication by three on
\(\widehat{\mathbf Z}^2\) imply that its fixed submodule
is zero. Hence \(T_e\subseteq E_1'\).

The actual map \(q_3:E_1\to\mathrm{SL}_2(\mathbf F_3)\)
is well defined because \(N_1\subseteq\ker q\), and it sends
\(E_1'\) onto the quaternion group \(Q_8\) of A.105.
Put, now over \(\mathbf F_3\),
\[
 \begin{gathered}
 X=\begin{pmatrix}1&2\\0&1\end{pmatrix},\quad
 Y=\begin{pmatrix}1&0\\1&1\end{pmatrix},\quad
 C=\begin{pmatrix}0&2\\1&2\end{pmatrix},\quad
 S=\begin{pmatrix}0&2\\1&0\end{pmatrix},\\
 D_e=\operatorname{diag}(1,e),\qquad
 \mathcal J_e=D_eCD_e^{-1}=(-1)^{\varepsilon}Y^{\varepsilon}C.
 \end{gathered}
\tag{A.111.13}
\]
The signed cyclic covariance is honest conjugation on the
derived subgroup: the extra exponent-parity character vanishes
there. Thus \(q_3(T_e)\subseteq C_{Q_8}(\mathcal J_e)\).
The centralizer is \(\{I,-I\}\). For \(e=1\) this is A.105's
\(C_{\mathrm{SL}_2(\mathbf F_3)}(C)=\langle C,-I\rangle\)
intersected with \(Q_8\); for \(e=-1\), conjugation by \(D_e\)
preserves the characteristic derived subgroup \(Q_8\) and gives
the same conclusion. This proves the target in (A.111.7).

Write \(J_e^{\alpha}=sj_e\), where
\[
 j_e=\begin{cases}\phi A_r,&\varepsilon=0,\\
 \iota(A_r),&\varepsilon=1,
 \end{cases}
 \qquad \iota=\operatorname{Ad}(s)|_F.
\tag{A.111.14}
\]
Both \(j_e\) lie in \(F'\). The actions of \(\alpha\),
\(\iota\), and inner conjugation by \(j_e\) on the signed
matrix image are matrix conjugations. They fix scalar matrices.
Since \(\kappa_e(t)=j_e^{-1}\iota(\alpha(t))j_e\), the map
\(\chi_e\) is equivariant with trivial action on its target.
It consequently induces the asserted map on pointed sets.
Here \(H^1(C_2,\{I,-I\})=\{I,-I\}\): every target element
is a cocycle and every target coboundary is \(I\).

\emph{A.111.4. Compute the image using the original frames.}
Use \(\nu\in\mathbf F_3\) from the ORIGINAL \(s\)-frame
coordinate in A.105. Its proved finite identities, together with
the actual full-GT conclusion of A.103, give
\[
 \begin{gathered}
 U=q_3(u)=eY^{-\nu}X^{\nu},\qquad
 \Phi=q_3(\phi)=Y^{\nu}X^{-\nu},\qquad
 A=q_3(A_r)=\epsilon_A X^{\nu}C^{-\nu},\\
 b=U_{12},\quad
 H=\begin{pmatrix}1&-b\\0&e\end{pmatrix},\quad
 K=H\operatorname{diag}(1,-1)H^{-1}
        =\begin{pmatrix}1&\nu\\0&-1\end{pmatrix},\qquad
 q_3\alpha=\operatorname{Ad}(K)q_3.
 \end{gathered}
\tag{A.111.15}
\]
There is no new choice of \(\Phi,A,U\) in these formulas.

A.108 computes, for every pure frame \(h\),
\[
 h_{\rm ab}=(t,-t),\quad t=(k+\varepsilon)/3,
 \quad t\bmod2=\varepsilon,
 \qquad Q\mathcal J_eQ^{-1}=HCH^{-1},\quad Q=q_3(h).
\tag{A.111.16}
\]
The quotient abelianization of the matrix group sends both
\(X,Y\) to the same generator of \(C_3\). Therefore the
displayed sum-zero abelianization implies \(Q\in Q_8\).
For clarity, the last identity in (A.111.16) retains the actual
cyclic sign: from
\(\Psi_1(r)=u^{-1}y^{-(k+\varepsilon)}R_e=hR_eh^{-1}\)
and signed covariance, its intermediate form is
\[
 Q\mathcal J_eQ^{-1}
   =(-1)^{\varepsilon-t}U^{-1}Y^{-k}C
   =U^{-1}Y^{-k}C=HCH^{-1}.
\]
Thus no map sending the full group element \(r\) directly to
the signed matrix \(C\) has been assumed.

Finally \(\iota(u^{-1})=u\) by canonical duality, so the word
in (A.111.5), rewritten entirely inside the pure kernel, is
\[
 w_h=j_e^{-1}\iota(\alpha(h)^{-1})u h.
\]
Hence its signed matrix is
\[
 \begin{split}
 W&=q_3(w_h)=q_3(j_e)^{-1}SKQ^{-1}K^{-1}S^{-1}UQ,\\
 q_3(j_e)&=
 \begin{cases}\Phi A,&e=1,\\ SAS^{-1},&e=-1.
 \end{cases}
 \end{split}
\tag{A.111.17}
\]
Here is the complete small calculation. For each pair \((\nu,e)\)
the equation in (A.111.16) has exactly the two solutions \(\pm Q\)
in \(Q_8\): their quotient centralizes \(\mathcal J_e\), whose
quaternion centralizer was just computed. Representatives and the
result of (A.111.17) with \(\epsilon_A=1\) are
\[
\begin{array}{c|c|c|c}
 \nu&e&Q\text{ up to sign}&W\ (\epsilon_A=1)\\\hline
 0& 1&\left(\begin{smallmatrix}1&0\\0&1\end{smallmatrix}\right)& I\\
 0&-1&\left(\begin{smallmatrix}1&0\\0&1\end{smallmatrix}\right)&-I\\
 1& 1&\left(\begin{smallmatrix}1&1\\1&2\end{smallmatrix}\right)& I\\
 1&-1&\left(\begin{smallmatrix}0&1\\2&0\end{smallmatrix}\right)&-I\\
 2& 1&\left(\begin{smallmatrix}0&1\\2&0\end{smallmatrix}\right)& I\\
 2&-1&\left(\begin{smallmatrix}1&2\\2&2\end{smallmatrix}\right)&-I
\end{array}
\tag{A.111.18}
\]
Each entry follows by substitution of (A.111.15) in (A.111.17).
Changing \(Q\) to \(-Q\) leaves \(W\) unchanged, while changing
\(\epsilon_A\) changes \(W\) by that scalar. This proves
\(W=\epsilon_A eI\) for all twenty-four possibilities.
The exact finite program and its complete output are retained as
\path{code/check_coherent_C2_CM34.py} and
\path{records/coherent_C2_CM34_certificate.json}; it enumerates all
twenty-four matrices of \(\mathrm{SL}_2(\mathbf F_3)\), all eight
elements of \(Q_8\), and all twenty-four admissible signed cases.
It exited successfully; the finite computation verifies this
comparison and is not used to infer profinite existence.

A.105 gives the original, already source-bound identity
\(\epsilon_Ae=(-1)^{\Delta(\alpha)}\). Since target cohomology
has no nontrivial coboundaries, (A.111.17)--(A.111.18) prove
(A.111.8) for the class itself. The final limitations in the
statement follow: a map of pointed sets need not have trivial
kernel, and neither the value of its present argument nor
that kernel has been decided here.
\end{proof}

This construction uses the quotient factor frame only through a
well-defined cohomology class. It does not restore uniqueness of
quotient frames or extend the arbitrary-frame formula of
LM7\_PUSH~E.65 to all unframed automorphisms of \(B_1\).
